\documentclass[11pt]{article}
\usepackage[margin=1in]{geometry}
\usepackage{amsmath}
\usepackage[hidelinks]{hyperref}

\title{A counterexample to Conjecture 00000003466}
\author{}
\date{}

\begin{document}
\maketitle

\noindent
The conjecture claims that the triangle core preserves the chromatic number
of every chordal graph and that the chromatic number of a general graph falls
by at most one under core contraction. The chordal graph $K_3$ loses two
colors when its unique triangle is contracted.

\section*{The contraction and the counterexample}

For a finite simple graph $G$ and a triangle $T$, let $q$ be the surjection
that identifies the three vertices of $T$ and leaves every other vertex in
its own class. Distinct quotient vertices are adjacent exactly when an edge
of $G$ joins the corresponding classes. This is also the simple graph
obtained by successively contracting two edges of $T$.

Take the complete graph $G=K_3$. It is chordal because it has no cycle of
length at least four. Its three vertices are pairwise adjacent, so a proper
coloring uses at least three colors, and assigning a distinct color to each
vertex gives $\chi(G)=3$. There is exactly one triangle. Contracting it leaves
one vertex with no edges, so the contracted graph has chromatic number one.
There are no triangles left and no contraction-order choices. Thus the
ordinary chromatic number is not preserved even in the chordal case, and
the drop $3-1=2$ exceeds the proposed bound for general graphs.

The example can also retain an edge after contraction. Attach a pendant
vertex to one vertex of $K_3$. All cycles still lie in the triangle, so the
new graph is chordal. Its chromatic number is three because the pendant
vertex can reuse a color of a nonneighbor in the triangle. Contracting the
unique triangle leaves $K_2$, whose chromatic number is two. This gives a
second failure of the chordal invariance claim.

\section*{Coloring constraints and Lean verification}

The quotient keeps the adjacency constraints between distinct classes.
Within the contracted class, a single color variable cannot express the
three pairwise inequalities on the original triangle. Retaining parallel
edges between distinct classes leaves the vertex-coloring rule unchanged.
Keeping an internal edge as a loop instead makes the quotient impossible to
color properly. An extra label encoding three independent color variables
would specify a coloring problem on the labeled classes in addition to the
ordinary chromatic number of the underlying graph.

The Lean 4 theorem \texttt{C3466.counterexample} gives the $K_3$ witness. It
checks chordality using cyclic orders and chords, verifies that the quotient
merges exactly the selected vertices, and proves both exact chromatic numbers.
It negates the claimed preservation and the bound on the loss for that
contraction. Run \texttt{lake build} in the \texttt{lean} directory to verify it.

\end{document}
